\usepackage[a4paper,margin=2.4cm,headheight=16pt]{geometry}
\usepackage{amsmath,amssymb,bm,mathtools,amsthm}
\usepackage{booktabs,array,longtable,tabularx}
\usepackage{graphicx,tikz,pgfplots}
\usetikzlibrary{arrows.meta,calc}
\pgfplotsset{compat=1.18}
\usepackage{fancyhdr,enumitem,float,placeins,microtype,needspace}
\usepackage[most]{tcolorbox}
\usepackage[colorlinks,linkcolor=blue!45!black,citecolor=green!35!black,urlcolor=blue!55!black]{hyperref}
\usepackage{xurl}
\setlength{\emergencystretch}{3em}
\setlist{nosep,leftmargin=2em}
\raggedbottom
\pagestyle{fancy}
\fancyhf{}
\fancyhead[L]{\small\nouppercase{\leftmark}}
\fancyfoot[C]{\thepage}
\newcommand{\dd}{\mathop{}\!\mathrm d}
\newcommand{\R}{\mathbb R}
\newcommand{\K}{\mathcal K}
\newcommand{\norm}[1]{\left\lVert #1\right\rVert}
\DeclareMathOperator{\diag}{diag}
\DeclareMathOperator{\tr}{tr}
\DeclareMathOperator{\nnz}{nnz}
\DeclareMathOperator{\curl}{curl}
\DeclareMathOperator{\diverg}{div}
\pdfstringdefDisableCommands{\def\textsuperscript#1{#1}}
\newtcolorbox{chaptersummary}{colback=blue!3,colframe=blue!45!black,title=本章总结,fonttitle=\bfseries,arc=1mm}
\newtcolorbox{authornote}{colback=yellow!7,colframe=orange!70!black,title=潜在认知跳跃,fonttitle=\bfseries,arc=1mm}
\newtcolorbox{askai}{colback=violet!3,colframe=violet!50!black,title=你应该问 AI,fonttitle=\bfseries,arc=1mm}
\newcommand{\learninggoals}[1]{\paragraph{学习目标}#1}
\newcommand{\studyroute}[2]{\paragraph{展开与检索}#1\par\textbf{关键词：}#2}
\newcommand{\conceptbridge}[3]{\begin{authornote}\textbf{潜在疑问：}#1\par\textbf{作者补充：}#2\par\textbf{核心概念联系：}#3\end{authornote}}
\newcommand{\aiexercise}[2]{\begin{askai}\textbf{AI 展开问题：}#1\par\textbf{进一步展开：}#2\end{askai}}
\newcommand{\costnote}[1]{\paragraph{工程成本与适用边界}#1}
\newcommand{\crossref}[1]{\paragraph{知识接口}#1}
\newcommand{\coderef}[1]{\paragraph{计算案例}位于本卷 \texttt{code\_examples/}：\par{\small\raggedright\nolinkurl{#1}\par}}
